\section{Implementation contract}
\label{sec:implementation}

The numerical implementation is organized into four layers.

\begin{enumerate}
\item \textbf{Electronic model.} Load and canonicalize the spinor Wannier Hamiltonian, time-reversal metadata, exchange-field decomposition, magnetic projectors, and local frames.
\item \textbf{Mixed electronic response.} Stream screened DFPT perturbations, build wrapped Green functions, accumulate the complex retarded loop, and complete q pairs.
\item \textbf{Mode projection.} Convert Cartesian kernels to phonon zero-point modes and then to external magnon quasiparticles.
\item \textbf{Validation and output.} Store HDF5 metadata, source hashes, restart status, decomposition residuals, q-symmetry tests, subspace sensitivity, and convergence tables.
\end{enumerate}

The reference matrix contraction for a batch of torque vertices $T_t$ and displacement perturbations $g_p$ uses
\begin{equation}
A_p=G_{\kk+\qq}g_pG_\kk,
\qquad
L_{tp}=\tr[T_tA_p],
\label{eq:batch-contraction}
\end{equation}
without materializing a tensor carrying both vertex labels and two matrix indices. q pairs are the primary parallel unit, while deterministic serial reductions remain available for validation.

All input files declare Bloch and Fourier gauges, basis ordering, time-reversal sewing convention, magnetic orbital masks, site projection, spin coordinate, and DFPT normalization. Missing metadata is an error rather than an invitation to infer conventions from array shapes.
